\documentclass[tikz,border=8pt]{standalone}
\usepackage{tikz}
\usepackage{amsmath}
\usepackage{pgfplots}
\pgfplotsset{compat=1.18}
\usepackage{ctex}
\usetikzlibrary{calc, positioning, arrows.meta, shapes, fit, backgrounds}

\begin{document}
\begin{tikzpicture}[
  >=Stealth,
  mrf node/.style={circle, draw=blue!70!black, fill=blue!10,
                   minimum size=12mm, font=\normalsize\bfseries, thick},
  fg var/.style={circle, draw=blue!70!black, fill=blue!10,
                 minimum size=10mm, font=\small\bfseries, thick},
  fg fac/.style={rectangle, draw=orange!80!black, fill=orange!15,
                 minimum size=8mm, font=\small, thick, rounded corners=2pt},
]

% ============================================================
% 左子图：马尔可夫随机场（无向图，4节点完全连接）
% ============================================================

% 节点位置（菱形布局）
\node[mrf node] (mA) at (0.0, 3.0)  {$A$};
\node[mrf node] (mB) at (2.5, 4.5)  {$B$};
\node[mrf node] (mC) at (5.0, 3.0)  {$C$};
\node[mrf node] (mD) at (2.5, 1.5)  {$D$};

% 无向边（完全图，共6条边）
\draw[blue!50!black, thick] (mA) -- (mB);
\draw[blue!50!black, thick] (mA) -- (mC);
\draw[blue!50!black, thick] (mA) -- (mD);
\draw[blue!50!black, thick] (mB) -- (mC);
\draw[blue!50!black, thick] (mB) -- (mD);
\draw[blue!50!black, thick] (mC) -- (mD);

% 高亮团 {A, B, C}（橙色填充背景）
\begin{scope}[on background layer]
  \node[draw=orange!70, fill=orange!20, thick, ellipse,
        fit=(mA)(mB)(mC), inner sep=2pt] {};
\end{scope}

% 势函数标注
\node[font=\scriptsize, orange!70!black] at (2.5, 5.5) {团 $\{A,B,C\}$：$\psi(A,B,C)$};

% MRF 子图标题
\node[font=\normalsize\bfseries] at (2.5, 6.5) {马尔可夫随机场（MRF）};

% MRF 公式
\node[font=\small,
      fill=white, draw=gray!50, thin, rounded corners=2pt, inner sep=4pt]
  at (2.5, 0.0) {$P(\mathbf{x})\propto \prod_C \psi_C(\mathbf{x}_C)$};

% ============================================================
% 右子图：因子图（二分图结构）
% ============================================================

% 变量节点（圆形）
\node[fg var] (fX1) at (8.0, 5.5) {$X_1$};
\node[fg var] (fX2) at (8.0, 3.5) {$X_2$};
\node[fg var] (fX3) at (8.0, 1.5) {$X_3$};

% 因子节点（方形）
\node[fg fac] (f12)  at (11.0, 5.5) {$f_{12}$};
\node[fg fac] (f23)  at (11.0, 3.5) {$f_{23}$};
\node[fg fac] (f123) at (11.0, 1.5) {$f_{123}$};

% 边（变量节点与因子节点之间）
\draw[blue!50!black, thick] (fX1) -- (f12);
\draw[blue!50!black, thick] (fX2) -- (f12);
\draw[blue!50!black, thick] (fX2) -- (f23);
\draw[blue!50!black, thick] (fX3) -- (f23);
\draw[blue!50!black, thick] (fX1) -- (f123);
\draw[blue!50!black, thick] (fX2) -- (f123);
\draw[blue!50!black, thick] (fX3) -- (f123);

% 因子图标题
\node[font=\normalsize\bfseries] at (9.5, 6.8) {因子图（Factor Graph）};

% 因子图图例
\node[fg var, minimum size=7mm, font=\scriptsize] at (8.2, -0.3) {};
\node[font=\scriptsize] at (9.0, -0.3) {变量节点};
\node[fg fac, minimum size=7mm, font=\scriptsize] at (10.5, -0.3) {};
\node[font=\scriptsize] at (11.5, -0.3) {因子节点};

% 因子图公式
\node[font=\small,
      fill=white, draw=gray!50, thin, rounded corners=2pt, inner sep=4pt]
  at (9.5, -1.2) {$P(\mathbf{x})\propto f_{12}(x_1,x_2)\,f_{23}(x_2,x_3)\,f_{123}(x_1,x_2,x_3)$};

% 分隔线
\draw[gray!40, thin, dashed] (6.5, -1.8) -- (6.5, 7.2);

% 整体大标题
\node[font=\large\bfseries] at (7.0, 8.0) {概率图模型：MRF 与因子图};

\end{tikzpicture}
\end{document}
